% !TEX root = ../main.tex
%==============================================================
\chapter{Nền tảng lý thuyết}
\label{ch:ly-thuyet}

Chương này trình bày cơ sở lý thuyết của hệ thống. Trước hết ta xây dựng các đại
lượng trắc quang (radiometry) và phương trình kết xuất mô tả sự cân bằng năng
lượng ánh sáng. Tiếp theo là phương pháp Monte Carlo dùng để ước lượng tích phân
trong phương trình đó, kèm các kỹ thuật giảm phương sai. Phần sau trình bày mô
hình hàm phản xạ (BRDF) cho ba loại vật liệu được hiện thực, mô hình camera thấu
kính mỏng, bài toán tìm giao tia--vật thể, và cuối cùng là mô hình tính toán song
song trên CUDA. Nội dung được tham khảo chủ yếu từ loạt sách \emph{Ray Tracing in
One Weekend}~\cite{shirley2020raytracing, shirley2020restoflife} và sách
\emph{Physically Based Rendering}~\cite{pharr2016pbrt}.

\section{Đại lượng trắc quang và phương trình kết xuất}

\subsection{Góc khối và bức xạ}
Để định lượng ánh sáng một cách độc lập với hình học, ta dùng khái niệm \emph{góc
khối} (solid angle). Một vùng trên mặt cầu đơn vị chắn bởi một hướng nhìn có độ
lớn góc khối $d\omega$ (đơn vị steradian). Trong hệ tọa độ cầu, vi phân góc khối
quanh hướng $(\theta,\phi)$ là
\begin{equation}
d\omega = \sin\theta\, d\theta\, d\phi .
\end{equation}

Đại lượng trung tâm của kết xuất dựa trên vật lý là \emph{bức xạ} (radiance)
$L$: công suất ánh sáng trên một đơn vị diện tích vuông góc, trên một đơn vị góc
khối, ký hiệu $L(\mathbf{p}, \omega)$ tại điểm $\mathbf{p}$ theo hướng $\omega$
\cite{pharr2016pbrt}. Bức xạ là đại lượng được \emph{bảo toàn dọc theo tia} trong
môi trường chân không và là giá trị mà mỗi tia trong path tracer mang theo. Từ
bức xạ tới, \emph{độ rọi} (irradiance) $E$ trên một bề mặt có pháp tuyến
$\mathbf{n}$ thu được bằng cách lấy tích phân theo bán cầu, có trọng số cosine:
\begin{equation}
\label{eq:irradiance}
E(\mathbf{p}) = \int_{\Omega} L_i(\mathbf{p}, \omega_i)\,
                (\omega_i \cdot \mathbf{n})\, d\omega_i .
\end{equation}
Thừa số cosine $(\omega_i\cdot\mathbf{n})=\cos\theta_i$ phản ánh việc một chùm
sáng tới xiên trải năng lượng trên diện tích lớn hơn nên đóng góp ít hơn.

\begin{figure}[H]
\centering
\begin{tikzpicture}[scale=1.7,>=Stealth]
  % surface
  \fill[gray!12] (-1.7,-0.35) rectangle (1.7,0);
  \draw[thick] (-1.7,0) -- (1.7,0);
  \foreach \x in {-1.5,-1.2,...,1.5}
    \draw[gray] (\x,0) -- (\x-0.18,-0.18);
  % hemisphere
  \draw[dashed,gray!70] (-1.4,0) arc (180:0:1.4);
  % point
  \coordinate (p) at (0,0);
  % normal
  \draw[->,thick] (p) -- (0,1.5) node[above]{$\mathbf{n}$};
  % incoming ray
  \draw[->,thick,blue!70!black] (-1.19,1.19) -- (p);
  \node[blue!70!black] at (-1.0,1.32) {$\omega_i$};
  % outgoing ray
  \draw[->,thick,red!75!black] (p) -- (1.19,1.19);
  \node[red!75!black] at (1.0,1.32) {$\omega_o$};
  % theta_i angle
  \draw (0,0.55) arc (90:135:0.55);
  \node at (-0.18,0.72) {\scriptsize$\theta_i$};
  % solid angle patch
  \draw[fill=blue!15,draw=blue!50] (-0.92,0.92) ++(-0.12,0.05)
        arc (135:120:1.4) -- ++(0.1,0.13) arc (120:135:1.55) -- cycle;
  \node[blue!60!black] at (-1.05,0.62) {\scriptsize$d\omega_i$};
  \fill (p) circle (1.1pt);
  \node[below=2pt] at (p) {$\mathbf{p}$};
\end{tikzpicture}
\caption{Hình học tại một điểm bề mặt: bức xạ tới theo hướng $\omega_i$ trong góc
khối vi phân $d\omega_i$, bức xạ phản xạ theo hướng $\omega_o$, với pháp tuyến
$\mathbf{n}$ và góc tới $\theta_i$.}
\label{fig:hemisphere}
\end{figure}

\subsection{Hàm phản xạ lưỡng hướng (BRDF)}
Cách một bề mặt phản xạ ánh sáng được mô tả bởi \emph{hàm phân bố phản xạ lưỡng
hướng} (Bidirectional Reflectance Distribution Function -- BRDF), định nghĩa là
tỉ số giữa bức xạ phản xạ theo hướng $\omega_o$ và độ rọi vi phân tới từ hướng
$\omega_i$:
\begin{equation}
f_r(\mathbf{p}, \omega_i, \omega_o) =
   \frac{dL_o(\mathbf{p}, \omega_o)}{L_i(\mathbf{p}, \omega_i)\,
   (\omega_i \cdot \mathbf{n})\, d\omega_i} .
\end{equation}
Một BRDF hợp lệ về mặt vật lý phải thỏa mãn hai tính chất: \emph{tính đối xứng
Helmholtz} $f_r(\omega_i,\omega_o)=f_r(\omega_o,\omega_i)$, và \emph{bảo toàn
năng lượng} -- tổng năng lượng phản xạ không vượt quá năng lượng tới
\cite{pharr2016pbrt}.

\subsection{Phương trình kết xuất}
Kết hợp BRDF với độ rọi tới, bức xạ rời khỏi điểm $\mathbf{p}$ theo hướng
$\omega_o$ được cho bởi \emph{phương trình kết xuất} do Kajiya đề
xuất~\cite{kajiya1986rendering}:
\begin{equation}
\label{eq:rendering}
L_o(\mathbf{p}, \omega_o) = L_e(\mathbf{p}, \omega_o)
  + \int_{\Omega} f_r(\mathbf{p}, \omega_i, \omega_o)\,
    L_i(\mathbf{p}, \omega_i)\, (\omega_i \cdot \mathbf{n})\, d\omega_i,
\end{equation}
trong đó $L_e$ là bức xạ tự phát (vật liệu phát sáng), số hạng tích phân là bức
xạ phản xạ. Đây là phương trình \emph{Fredholm loại hai}: ẩn $L$ xuất hiện ở cả
hai vế vì bức xạ tới $L_i(\mathbf{p},\omega_i)$ tại $\mathbf{p}$ chính là bức xạ
đi ra $L_o(\mathbf{p}', -\omega_i)$ của điểm $\mathbf{p}'$ mà tia $(\mathbf{p},
\omega_i)$ chạm tới. Tính đệ quy này khiến lời giải giải tích là bất khả thi với
cảnh tổng quát, do đó cần phương pháp số.

\section{Ước lượng Monte Carlo}

\subsection{Bộ ước lượng cơ bản}
Phương pháp Monte Carlo xấp xỉ một tích phân xác định bằng kỳ vọng thống kê. Với
tích phân $I = \int g(x)\, dx$, nếu lấy $N$ mẫu độc lập $x_k$ theo mật độ xác
suất $p(x)$ (với $p(x)>0$ ở những nơi $g(x)\neq0$), \emph{bộ ước lượng Monte
Carlo} là
\begin{equation}
\label{eq:montecarlo}
\langle I \rangle = \frac{1}{N}\sum_{k=1}^{N} \frac{g(x_k)}{p(x_k)} .
\end{equation}
Bộ ước lượng này \emph{không chệch} (unbiased): $\mathbb{E}[\langle I\rangle]=I$
với mọi $N$. Phương sai của nó giảm tỉ lệ $1/N$, nên \emph{sai số chuẩn} (độ lệch
chuẩn) giảm theo $O(1/\sqrt{N})$. Hệ quả thực tiễn: muốn giảm một nửa nhiễu cần
gấp bốn lần số mẫu -- đây là nguồn gốc của hiện tượng \emph{nhiễu hạt} (noise)
đặc trưng và chi phí tính toán lớn của path tracing~\cite{shirley2020restoflife}.

\subsection{Lấy mẫu theo tầm quan trọng}
Phương sai phụ thuộc mạnh vào mức độ phù hợp giữa $p(x)$ và $g(x)$. Nếu chọn
$p(x)$ tỉ lệ với $g(x)$ thì tỉ số $g/p$ gần như hằng số và phương sai giảm mạnh;
đây là ý tưởng \emph{lấy mẫu theo tầm quan trọng} (importance
sampling)~\cite{veach1997robust}. Trong kết xuất, do tích phân
\eqref{eq:rendering} chứa thừa số cosine và BRDF, ta thường lấy mẫu hướng
$\omega_i$ theo phân bố cosine hoặc theo hình dạng của BRDF thay vì đều trên bán
cầu. Hệ thống hiện tại lấy mẫu khuếch tán theo phân bố cosine (Mục~\ref{sec:vatlieu});
việc lấy mẫu trực tiếp nguồn sáng (Next Event Estimation) là hướng mở rộng nêu ở
Chương~\ref{ch:ket-luan}.

\subsection{Công thức tích phân đường và hệ số suy giảm}
Khai triển đệ quy phương trình \eqref{eq:rendering} cho \emph{công thức tích phân
đường} (path integral): bức xạ tới camera là tổng đóng góp của mọi đường ánh sáng
nối nguồn sáng tới mắt qua một số lần phản xạ. Path tracing ước lượng tổng này
bằng cách, tại mỗi điểm va chạm, lấy \emph{một} hướng tán xạ ngẫu nhiên và lần
theo tia đó -- tạo thành một đường đi ngẫu nhiên. Đóng góp của đường được tích
lũy qua \emph{hệ số suy giảm} (throughput) $\beta$, cập nhật nhân dồn sau mỗi va
chạm:
\begin{equation}
\label{eq:throughput}
\beta_n = \beta_{n-1}\,
          \frac{f_r(\mathbf{p}_n,\omega_i,\omega_o)\,(\omega_i\cdot\mathbf{n})}
               {p(\omega_i)},
\qquad \beta_0 = 1 .
\end{equation}
Khi tia thoát khỏi cảnh và gặp nguồn sáng nền, bức xạ thu được bằng tích của bức
xạ nền với $\beta$. Với cách lấy mẫu khuếch tán theo cosine, thừa số cosine và
mẫu số $p(\omega_i)$ triệt tiêu nhau, khiến việc cập nhật $\beta$ rút gọn về phép
nhân với albedo -- đúng như hiện thực trong mã nguồn (Chương~\ref{ch:trien-khai}).

\begin{figure}[H]
\centering
\begin{tikzpicture}[scale=1.0,>=Stealth,
   surf/.style={line width=1.2pt,gray!70}]
  % camera
  \node[draw,fill=gray!10,minimum width=0.7cm,minimum height=0.5cm] (cam) at (0,0) {};
  \node[below=1pt of cam,font=\small] {Camera};
  % bounce points
  \coordinate (b1) at (2.3,0.9);
  \coordinate (b2) at (4.4,-0.4);
  \coordinate (b3) at (6.2,1.1);
  % surfaces (short segments) with normals
  \draw[surf] ($(b1)+(-0.45,-0.25)$) -- ($(b1)+(0.45,0.25)$);
  \draw[surf] ($(b2)+(-0.5,0.1)$) -- ($(b2)+(0.5,-0.1)$);
  \draw[surf] ($(b3)+(-0.4,-0.3)$) -- ($(b3)+(0.4,0.3)$);
  % light source
  \node[star,star points=8,star point ratio=0.6,fill=yellow!80!orange,
        draw=orange,minimum size=0.9cm] (L) at (7.8,2.3) {};
  \node[above=0pt of L,font=\small] {Nguồn sáng};
  % path segments
  \draw[->,thick,blue!70!black] (cam.east) -- (b1) node[midway,above,font=\scriptsize]{$\beta_0$};
  \draw[->,thick,blue!70!black] (b1) -- (b2) node[midway,above,font=\scriptsize]{$\beta_1$};
  \draw[->,thick,blue!70!black] (b2) -- (b3) node[midway,below,font=\scriptsize]{$\beta_2$};
  \draw[->,thick,blue!70!black] (b3) -- (L) node[midway,above left,font=\scriptsize]{$\beta_3$};
  % bounce dots
  \foreach \b in {b1,b2,b3} \fill (\b) circle (1.6pt);
  \node[right=1pt,font=\scriptsize] at (b1) {$\mathbf{p}_1$};
  \node[right=1pt,font=\scriptsize] at (b2) {$\mathbf{p}_2$};
  \node[right=1pt,font=\scriptsize] at (b3) {$\mathbf{p}_3$};
\end{tikzpicture}
\caption{Một đường đi ngẫu nhiên trong Path Tracing: tia xuất phát từ camera, tán
xạ qua các điểm va chạm $\mathbf{p}_1,\mathbf{p}_2,\dots$ và kết thúc tại nguồn
sáng; hệ số suy giảm $\beta$ được nhân dồn dọc theo đường.}
\label{fig:pathtrace}
\end{figure}

\subsection{Kết thúc đường đi và Russian Roulette}
Vì một đường có thể phản xạ vô hạn, cần điều kiện dừng. Cách đơn giản là giới hạn
độ sâu tối đa, song điều này tạo ra sai lệch (bias) do bỏ qua các đường dài. Một
giải pháp không chệch là \emph{Russian Roulette}~\cite{pharr2016pbrt}: tại mỗi
bước, kết thúc đường với xác suất $q$ và chia phần còn lại cho $1-q$ để giữ kỳ
vọng không đổi. Hệ thống áp dụng giới hạn độ sâu (tối đa 50 lần dội) kết hợp dừng
sớm khi $\beta$ suy giảm gần bằng 0 -- một xấp xỉ thực dụng cho chế độ thời gian
thực.

\section{Lấy mẫu điểm ảnh và tích lũy theo thời gian}

\subsection{Lý thuyết lấy mẫu và hiện tượng răng cưa}
Quá trình tạo ảnh số thực chất là \emph{lấy mẫu} (sampling) một hàm ảnh liên tục
$f(x,y)$ -- giá trị bức xạ tại mọi điểm trên mặt phẳng ảnh -- tại một lưới điểm
ảnh rời rạc, rồi \emph{tái dựng} (reconstruction) ảnh hiển thị từ các mẫu đó. Cơ
sở lý thuyết của quá trình này là phân tích Fourier~\cite{pharr2016pbrt}.

Mỗi hàm $f$ có một \emph{phổ tần số} $F = \mathcal{F}\{f\}$ thu được qua biến đổi
Fourier. Việc lấy mẫu đều với bước $T$ tương đương nhân $f$ với một chuỗi xung
Dirac (Dirac comb) khoảng cách $T$. Theo định lý tích chập, trong miền tần số
phép nhân này trở thành tích chập của phổ $F$ với một chuỗi xung khoảng cách
$1/T$, tạo ra \emph{vô số bản sao} của phổ gốc lặp lại theo chu kỳ $1/T$:
\begin{equation}
\mathcal{F}\Big\{ f \cdot \textstyle\sum_n \delta(x - nT)\Big\}
   = F * \frac{1}{T}\sum_k \delta\!\Big(\nu - \frac{k}{T}\Big).
\end{equation}
Để khôi phục $f$, ta lọc lấy bản sao trung tâm bằng một bộ lọc thông thấp lý
tưởng. Điều này khả thi \emph{nếu và chỉ nếu} các bản sao không chồng lấn, tức
hàm phải \emph{giới hạn dải tần} (band-limited) với tần số lớn nhất $\nu_{\max}$,
và tần số lấy mẫu thỏa
\begin{equation}
\label{eq:nyquist}
\frac{1}{T} \;>\; 2\,\nu_{\max} .
\end{equation}
Bất đẳng thức~\eqref{eq:nyquist} chính là \emph{định lý lấy mẫu
Nyquist--Shannon}; ngưỡng $2\nu_{\max}$ gọi là \emph{tần số Nyquist}.

Trong kết xuất, điều kiện này gần như luôn bị vi phạm: biên hình học, bóng đổ
sắc nét hay hoa văn texture tạo ra các \emph{bước nhảy} có phổ trải tới vô hạn
nên ảnh \emph{không} giới hạn dải tần. Khi đó các bản sao phổ chồng lấn, năng
lượng tần số cao ``giả dạng'' thành tần số thấp sau khi tái dựng, sinh ra các
sai lệch nhìn thấy được gọi là \emph{răng cưa} (aliasing): biên xiên bị bậc
thang, hoa văn nhỏ biến thành vân moiré.

\begin{figure}[H]
\centering
\begin{tikzpicture}[scale=1.0,>=Stealth]
  % ---- top: adequate sampling ----
  \begin{scope}[yshift=2.2cm]
    \draw[->] (-3.2,0) -- (3.2,0) node[right,font=\scriptsize]{$\nu$};
    \foreach \c in {-2,0,2} {
      \draw[blue!60!black,fill=blue!12] (\c-0.7,0) -- (\c,0.8) -- (\c+0.7,0) -- cycle;
    }
    \node[font=\scriptsize] at (0,1.0) {phổ gốc};
    \node[font=\scriptsize,gray] at (2.1,0.95) {bản sao};
    \node[left,font=\scriptsize] at (-3.2,0.55) {(a)};
  \end{scope}
  % ---- bottom: undersampling -> overlap ----
  \begin{scope}
    \draw[->] (-3.2,0) -- (3.2,0) node[right,font=\scriptsize]{$\nu$};
    \foreach \c in {-1.1,0,1.1} {
      \draw[red!60!black,fill=red!12] (\c-0.7,0) -- (\c,0.8) -- (\c+0.7,0) -- cycle;
    }
    % overlap markers
    \fill[red!45,opacity=0.6] (-0.55,0) -- (-0.55,0.2) -- (-0.4,0) -- cycle;
    \fill[red!45,opacity=0.6] (0.55,0) -- (0.55,0.2) -- (0.4,0) -- cycle;
    \node[font=\scriptsize,red!70!black] at (0.0,1.0) {vùng chồng lấn $\to$ răng cưa};
    \node[left,font=\scriptsize] at (-3.2,0.55) {(b)};
  \end{scope}
\end{tikzpicture}
\caption{Phổ tần số sau khi lấy mẫu gồm phổ gốc và các bản sao lặp lại. (a) Khi
thỏa điều kiện Nyquist~\eqref{eq:nyquist}, các bản sao tách rời nên khôi phục
được tín hiệu. (b) Khi lấy mẫu thiếu (hoặc tín hiệu không giới hạn dải tần), các
bản sao chồng lấn, gây hiện tượng răng cưa.}
\label{fig:nyquist}
\end{figure}

\subsection{Khử răng cưa bằng lấy mẫu ngẫu nhiên}
Có hai hướng chống răng cưa. Hướng thứ nhất là \emph{tiền lọc} (prefiltering) --
làm mượt hàm ảnh để giới hạn dải tần trước khi lấy mẫu -- nhưng nhìn chung không
khả thi vì hàm ảnh không có dạng giải tích. Hướng thứ hai, thực dụng hơn, là
\emph{tăng tần số lấy mẫu}: bắn nhiều tia trên mỗi điểm ảnh (\emph{siêu lấy mẫu},
supersampling) để đẩy tần số Nyquist lên cao, thu hẹp vùng bị răng cưa.

Điều quan trọng là path tracer dùng \emph{lấy mẫu ngẫu nhiên} (stochastic
sampling) thay vì lưới đều. Với mẫu đều, năng lượng tần số cao chồng lấn tạo ra
các vân răng cưa có cấu trúc; còn khi điểm lấy mẫu được làm \emph{nhiễu} ngẫu
nhiên, năng lượng đó bị phân tán thành \emph{nhiễu} (noise) không cấu trúc. Mắt
người ít nhạy với nhiễu hơn nhiều so với các vân răng cưa có quy luật, nên đây là
sự đánh đổi có lợi~\cite{pharr2016pbrt}.

Về mặt hình thức, giá trị một điểm ảnh là tích phân của hàm ảnh nhân với một
\emph{hàm lọc tái dựng} $r$ trên miền điểm ảnh, $I_{ij} = \int r(x,y)\,f(x,y)\,
dx\,dy$. Path tracer ước lượng tích phân này bằng Monte Carlo
\eqref{eq:montecarlo} với các mẫu được làm nhiễu ngẫu nhiên trong ô điểm ảnh.
Mỗi khung dịch điểm bắn tia một lượng $(\xi_u,\xi_v)\in[0,1)^2$,
\begin{equation}
s = \frac{i + \xi_u}{W}, \qquad t = 1 - \frac{j + \xi_v}{H},
\end{equation}
với $(i,j)$ là chỉ số điểm ảnh và $W\times H$ là độ phân giải (hệ thống dùng hàm
lọc hộp, tức trọng số đều trong ô). Trung bình của nhiều mẫu vừa khử răng cưa
biên, vừa hội tụ về giá trị tích phân đúng~\cite{shirley2020raytracing}.

\subsection{Tích lũy theo thời gian}
Trong chế độ tương tác, mỗi khung chỉ kịp lấy một mẫu/điểm ảnh nên ảnh rất nhiễu.
\emph{Tích lũy theo thời gian} (temporal accumulation) tận dụng việc camera đứng
yên: màu của các khung liên tiếp được cộng dồn rồi lấy trung bình,
\begin{equation}
C_{\text{out}} = \frac{1}{F}\sum_{f=1}^{F} C_f ,
\end{equation}
trong đó $F$ là số khung kể từ lần đặt lại gần nhất. Đây chính là bộ ước lượng
\eqref{eq:montecarlo} được tích lũy dần theo thời gian: càng nhiều khung, $F$ càng
lớn, ảnh càng hội tụ (giảm nhiễu). Khi người dùng di chuyển camera, $F$ được đặt
lại về 1 để loại bỏ các mẫu cũ không còn hợp lệ.

\subsection{Hiệu chỉnh gamma}
Tính toán ánh sáng diễn ra trong không gian màu \emph{tuyến tính}, nhưng màn hình
kỳ vọng giá trị đã mã hóa gamma. Trước khi hiển thị, màu được hiệu chỉnh bằng
$C' = C^{1/\gamma}$; hệ thống dùng $\gamma = 2$, tức lấy căn bậc hai từng kênh,
giúp vùng tối hiển thị đúng độ sáng cảm nhận~\cite{shirley2020raytracing}.

\section{Mô hình vật liệu}
\label{sec:vatlieu}
Mỗi vật liệu định nghĩa một hàm \emph{tán xạ}: cho tia tới, trả về tia tán xạ và
hệ số suy giảm (attenuation). Hệ thống hiện thực ba mô hình kinh điển của loạt
\emph{Ray Tracing in One Weekend}~\cite{shirley2020raytracing}.

\subsection{Vật liệu khuếch tán (Lambertian)}
Bề mặt khuếch tán lý tưởng có BRDF hằng số:
\begin{equation}
f_r = \frac{\rho}{\pi},
\end{equation}
với $\rho$ là albedo (hệ số phản xạ) và hệ số $\pi$ bảo đảm bảo toàn năng lượng
khi lấy tích phân theo bán cầu. Để lấy mẫu hiệu quả, ta sinh hướng tán xạ theo
\emph{phân bố cosine} bằng một mẹo hình học: cộng pháp tuyến với một điểm ngẫu
nhiên trên mặt cầu đơn vị,
\begin{equation}
\omega_i = \mathbf{n} + \mathbf{r}_{\text{unit}}, \qquad
\|\mathbf{r}_{\text{unit}}\| = 1 .
\end{equation}
Khi đó mật độ lấy mẫu $p(\omega_i)=\cos\theta_i/\pi$ khớp với thừa số cosine và
BRDF, làm phép cập nhật throughput \eqref{eq:throughput} rút gọn còn phép nhân
với albedo $\rho$~\cite{shirley2020restoflife}.

\subsection{Vật liệu kim loại (Metal)}
Kim loại phản xạ gương: tia phản xạ quanh pháp tuyến theo
\begin{equation}
\label{eq:reflect}
\mathbf{r} = \mathbf{d} - 2(\mathbf{d}\cdot\mathbf{n})\,\mathbf{n},
\end{equation}
với $\mathbf{d}$ là hướng tia tới đã chuẩn hóa. Để mô phỏng bề mặt nhám, hướng
phản xạ được nhiễu loạn bằng một vector ngẫu nhiên trong quả cầu bán kính
\emph{fuzz}: $\mathbf{r}' = \mathbf{r} + \text{fuzz}\cdot\mathbf{r}_{\text{rand}}$.
Giá trị fuzz càng lớn, phản chiếu càng mờ; tia tán xạ chui xuống dưới bề mặt sẽ
bị loại bỏ.

\subsection{Vật liệu điện môi (Dielectric)}
Vật liệu trong suốt (thủy tinh, nước) vừa phản xạ vừa khúc xạ. Phần khúc xạ tuân
theo \emph{định luật Snell}:
\begin{equation}
\label{eq:snell}
\eta \sin\theta_i = \eta' \sin\theta_t,
\end{equation}
với $\eta,\eta'$ là chiết suất hai môi trường, $\theta_i,\theta_t$ là góc tới và
góc khúc xạ. Khi $\frac{\eta}{\eta'}\sin\theta_i > 1$, phương trình
\eqref{eq:snell} vô nghiệm với $\sin\theta_t$: xảy ra \emph{phản xạ toàn phần},
tia buộc phải phản xạ theo \eqref{eq:reflect}.

Tỉ lệ năng lượng phản xạ so với khúc xạ được cho bởi \emph{phương trình Fresnel},
vốn phức tạp vì phụ thuộc phân cực. Trong đồ họa, ta dùng \emph{xấp xỉ Schlick}
chính xác cao và rẻ về tính toán~\cite{pharr2016pbrt}:
\begin{equation}
\label{eq:schlick}
R(\theta) = R_0 + (1 - R_0)\,(1 - \cos\theta)^5,
\qquad R_0 = \left(\frac{\eta - \eta'}{\eta + \eta'}\right)^2,
\end{equation}
trong đó $R_0$ là hệ số phản xạ ở góc tới vuông góc. Thay vì tính cả hai tia phản
xạ và khúc xạ rồi pha trộn, path tracer dùng cách \emph{lấy mẫu ngẫu nhiên}: sinh
một số $\xi\in[0,1)$ và chọn phản xạ nếu $\xi < R(\theta)$, ngược lại thì khúc xạ.
Theo nghĩa thống kê qua nhiều mẫu, tỉ lệ này tái tạo đúng hệ số Fresnel mà mỗi
tia chỉ phải lần theo một nhánh~\cite{shirley2020raytracing}.

\section{Mô hình camera thấu kính mỏng}

\subsection{Phương trình đo của camera}
Một cách hình thức, giá trị mà một điểm ảnh ghi nhận không phải là bức xạ của một
tia đơn lẻ, mà là một \emph{phép đo} (measurement) trên toàn bộ ánh sáng đi vào
hệ quang học. Theo PBRT~\cite{pharr2016pbrt}, đáp ứng $I_j$ của điểm ảnh thứ $j$
được cho bởi \emph{phương trình đo} (measurement equation):
\begin{equation}
\label{eq:measurement}
I_j = \int_{A_{\text{film}}} \int_{\Omega}
       W_e^{(j)}(\mathbf{p}, \omega)\, L_i(\mathbf{p}, \omega)\,
       |\cos\theta|\; d\omega\, dA(\mathbf{p}),
\end{equation}
trong đó tích phân lấy trên diện tích phim $A_{\text{film}}$ và bán cầu hướng
$\Omega$; $L_i$ là bức xạ tới, còn $W_e^{(j)}$ là \emph{hàm tầm quan trọng}
(importance) -- đặc trưng độ nhạy của cảm biến với tia $(\mathbf{p},\omega)$,
đóng vai trò đối ngẫu với bức xạ phát ra của một nguồn sáng. Nhiệm vụ của mô hình
camera là sinh các tia sao cho việc lấy mẫu Monte Carlo tích phân
\eqref{eq:measurement} là hợp lệ.

\subsection{Từ lỗ kim tới thấu kính mỏng}
Với camera \emph{lỗ kim} (pinhole) lý tưởng, khẩu độ thu về một điểm: hàm tầm
quan trọng $W_e$ suy biến thành một xung Dirac theo hướng, nên mỗi điểm trên phim
chỉ ứng với \emph{một} tia duy nhất và mọi vật thể đều sắc nét bất kể độ sâu.
Đây là lý do ảnh pinhole không có hiệu ứng nhòe của ống kính thật.

Để tái hiện \emph{độ sâu trường ảnh} (depth of field), ta cần một khẩu độ hữu
hạn. Mô hình \emph{thấu kính mỏng} (thin lens)~\cite{pharr2016pbrt} xấp xỉ hệ
quang học bằng một thấu kính lý tưởng độ dày bằng 0, tuân theo \emph{phương trình
thấu kính Gauss}:
\begin{equation}
\label{eq:thinlens}
\frac{1}{z'} - \frac{1}{z} = \frac{1}{f},
\end{equation}
với $z$ là khoảng cách từ vật tới thấu kính, $z'$ là khoảng cách ảnh, và $f$ là
tiêu cự. Phương trình \eqref{eq:thinlens} chỉ ra rằng các vật ở những khoảng cách
$z$ khác nhau hội tụ tại những mặt phẳng ảnh $z'$ khác nhau. Do cảm biến đặt cố
định, chỉ những vật nằm tại \emph{mặt phẳng lấy nét} (khoảng cách
\emph{focus distance}) mới hội tụ thành điểm sắc nét; vật ở khoảng cách khác trải
thành một đĩa gọi là \emph{vòng tròn nhòe} (circle of confusion), gây hiệu ứng mờ
nhòe. Bán kính vòng nhòe -- và do đó độ mạnh của hiệu ứng -- tỉ lệ với bán kính
khẩu độ.

\begin{figure}[H]
\centering
\begin{tikzpicture}[scale=1.0,>=Stealth]
  % sensor plane
  \draw[line width=1.4pt,gray!70] (0,-1.4) -- (0,1.4);
  \node[gray!70,font=\scriptsize] at (0,1.7) {Cảm biến};
  % lens
  \draw[blue!55!black,line width=1pt] (3,0) ellipse (0.18 and 1.2);
  \node[blue!55!black,font=\scriptsize] at (3,-1.5) {Thấu kính};
  \draw[<->,blue!55!black] (3.35,0) -- (3.35,1.2) node[midway,right,font=\scriptsize]{$R$};
  % focal plane
  \draw[dashed,green!45!black] (6.5,-1.5) -- (6.5,1.5);
  \node[green!45!black,font=\scriptsize,align=center] at (6.5,1.8) {Mặt phẳng\\lấy nét};
  % in-focus point (sharp): converges to single sensor point
  \coordinate (F) at (6.5,0.7);
  \fill[green!45!black] (F) circle (1.6pt);
  \coordinate (S) at (0,-0.55);
  \draw[green!50!black] (S) -- (3,0.95) -- (F);
  \draw[green!50!black] (S) -- (3,-0.95) -- (F);
  \fill[green!45!black] (S) circle (1.4pt);
  \node[green!45!black,font=\scriptsize,left] at (S) {nét};
  % out-of-focus point: spreads -> circle of confusion
  \coordinate (O) at (5.2,-0.6);
  \fill[red!70!black] (O) circle (1.6pt);
  \draw[red!65!black] (O) -- (3,0.95) -- (0,0.95);
  \draw[red!65!black] (O) -- (3,-0.95) -- (0,0.2);
  \draw[red!65!black,line width=1.3pt] (0,0.2) -- (0,0.95);
  \node[red!70!black,font=\scriptsize,right=1pt] at (O) {ngoài nét};
  \node[red!70!black,font=\scriptsize,align=left] at (-0.55,1.0) {vòng\\nhòe};
\end{tikzpicture}
\caption{Mô hình thấu kính mỏng. Điểm nằm trên mặt phẳng lấy nét (xanh) hội tụ
thành một điểm sắc nét trên cảm biến; điểm ngoài mặt phẳng lấy nét (đỏ) trải ra
thành vòng tròn nhòe, tạo hiệu ứng độ sâu trường ảnh.}
\label{fig:thinlens}
\end{figure}

\subsection{Sinh tia với khẩu độ hữu hạn}
Tích phân \eqref{eq:measurement} trên diện tích thấu kính được ước lượng bằng
Monte Carlo: thay vì một tia, ta lấy mẫu một điểm trên đĩa thấu kính bán kính
$R=\text{aperture}/2$. Quy trình sinh tia~\cite{pharr2016pbrt,shirley2020raytracing}
gồm ba bước:
\begin{enumerate}[noitemsep]
  \item Dựng tia chính đi từ tâm thấu kính qua điểm ảnh $(s,t)$.
  \item Xác định \emph{điểm lấy nét} là giao của tia chính với mặt phẳng lấy nét
        (cách camera một khoảng focus distance) -- điểm này luôn sắc nét.
  \item Lấy ngẫu nhiên một điểm trên đĩa thấu kính làm gốc tia mới, hướng tia mới
        đi qua đúng điểm lấy nét ở bước~2.
\end{enumerate}
Vì mọi tia lấy mẫu đều đi qua cùng điểm lấy nét, vật tại mặt phẳng lấy nét vẫn
sắc nét; còn vật ngoài mặt phẳng đó nhận các tia từ những vị trí thấu kính khác
nhau nên bị nhòe, đúng như mô tả của vòng tròn nhòe. Khi khẩu độ bằng 0, mô hình
thoái biến về camera lỗ kim.

\subsection{Hệ trục và nhòe chuyển động}
Trường nhìn dọc $\theta$ xác định chiều cao khung nhìn qua $h=\tan(\theta/2)$. Hệ
trục camera trực chuẩn $(\mathbf{u},\mathbf{v},\mathbf{w})$ được dựng từ vị trí
đặt máy (\emph{lookfrom}), điểm nhìn (\emph{lookat}) và vector hướng lên
(\emph{vup}):
\begin{equation}
\mathbf{w} = \frac{\text{lookfrom}-\text{lookat}}{\|\cdot\|},\quad
\mathbf{u} = \frac{\text{vup}\times\mathbf{w}}{\|\cdot\|},\quad
\mathbf{v} = \mathbf{w}\times\mathbf{u}.
\end{equation}
Để mô phỏng \emph{nhòe chuyển động} (motion blur), phương trình đo
\eqref{eq:measurement} được mở rộng thêm một tích phân theo \emph{thời gian phơi
sáng}. Hệ thống ước lượng tích phân này bằng cách gán cho mỗi tia một mốc thời
gian ngẫu nhiên trong khoảng $[t_0,t_1]$; vật thể chuyển động được nội suy vị trí
tuyến tính theo thời gian của tia, nên ảnh tích lũy thể hiện vệt mờ tự
nhiên~\cite{shirley2020raytracing}.

\section{Tìm giao tia--vật thể và cấu trúc tăng tốc}

\subsection{Giao tia--khối cầu}
Phép toán cốt lõi của mỗi bước dò tia là tìm giao điểm gần nhất giữa tia và vật
thể. Một tia tham số hóa $\mathbf{r}(t)=\mathbf{o}+t\mathbf{d}$ giao khối cầu tâm
$\mathbf{c}$, bán kính $r$ khi $\|\mathbf{r}(t)-\mathbf{c}\|^2 = r^2$. Khai triển
cho phương trình bậc hai $at^2 + bt + c = 0$ với
\begin{equation}
a = \mathbf{d}\cdot\mathbf{d},\quad
b = 2\,\mathbf{d}\cdot(\mathbf{o}-\mathbf{c}),\quad
c = (\mathbf{o}-\mathbf{c})\cdot(\mathbf{o}-\mathbf{c}) - r^2 .
\end{equation}
Biệt thức $\Delta=b^2-4ac$ âm nghĩa là không có giao; ngược lại, nghiệm nhỏ hơn
nằm trong khoảng $[t_{\min},t_{\max}]$ cho điểm va chạm gần nhất. Pháp tuyến tại
điểm giao $\mathbf{p}$ là $\mathbf{n}=(\mathbf{p}-\mathbf{c})/r$. Cận dưới
$t_{\min}$ dương nhỏ giúp tránh \emph{shadow acne} -- hiện tượng tia thứ cấp tự
giao lại bề mặt vừa xuất phát do sai số dấu phẩy động~\cite{shirley2020raytracing}.

\subsection{Cấu trúc tăng tốc}
Khi cảnh có $n$ vật thể, duyệt tuyến tính toàn bộ danh sách tốn $O(n)$ cho mỗi
tia, không khả thi với cảnh phức tạp. Các \emph{cấu trúc tăng tốc} như cây phân
cấp khối bao (Bounding Volume Hierarchy -- BVH) bao mỗi nhóm vật thể trong một
khối bao (thường là hộp căn trục AABB) và tổ chức thành cây nhị phân. Nhờ kiểm
tra giao với khối bao trước, ta loại bỏ sớm những nhánh tia không thể chạm, giảm
độ phức tạp trung bình xuống $O(\log n)$. Chất lượng cây phụ thuộc cách phân
hoạch; \emph{Surface Area Heuristic} (SAH) là tiêu chí phổ biến để chọn mặt phân
chia tối ưu~\cite{pharr2016pbrt}. Phiên bản hiện tại dùng duyệt tuyến tính; tích
hợp BVH là hướng tối ưu được bàn ở Chương~\ref{ch:ket-luan}.

\section{Kiến trúc lập trình song song CUDA}
CUDA là nền tảng tính toán song song đa dụng của NVIDIA, cho phép tận dụng hàng
nghìn nhân của GPU~\cite{nvidia2024cuda}. Path tracing có tính \emph{song song dữ
liệu} lý tưởng: mỗi điểm ảnh được tính độc lập, hoàn toàn phù hợp mô hình SIMT
(Single Instruction, Multiple Threads) của GPU.

Trong mô hình CUDA, công việc được tổ chức thành lưới (grid) gồm nhiều khối
(block), mỗi khối chứa nhiều luồng (thread) chạy đồng thời theo nhóm 32 luồng
(warp). Hàm chạy trên GPU gọi là \emph{kernel}; khi khởi chạy, mỗi luồng xác định
điểm ảnh phụ trách qua chỉ số toàn cục tính từ \texttt{blockIdx},
\texttt{blockDim} và \texttt{threadIdx}. Vì mỗi luồng cần một dãy số ngẫu nhiên
độc lập cho Monte Carlo, hệ thống dùng thư viện \texttt{cuRAND} với hạt giống
riêng cho từng điểm ảnh. Kích thước khối (ví dụ $8\times8=64$ luồng) ảnh hưởng
trực tiếp tới mức độ chiếm dụng (occupancy) và hiệu năng. Một điểm cần lưu ý là
sự \emph{phân kỳ luồng} (warp divergence): khi các tia trong cùng warp rẽ nhánh
khác nhau (ví dụ chạm vật liệu khác loại), GPU phải thực thi tuần tự các nhánh,
làm giảm hiệu suất -- đây là đặc thù quan trọng khi tối ưu path tracer trên GPU.
